\section{Monte Carlo evidence}\label{sec:mc}

Before turning to the data, we examine how EW-DRPL performs at sample sizes typical of
cross-country work, where the asymptotics behind Proposition~\ref{prop:regret} may be a poor guide.

\paragraph{Design.} Each simulated country has six correlated characteristics that mimic log income,
inequality, internet access, agricultural share, baseline temperature and knowledge sovereignty.
It is observed in $J=4$ spells, each with ten deciles. Capability is assigned with probability
$\mathrm{logit}^{-1}(1.5\,\text{income}+0.5\,\text{internet})$, so it is strongly confounded with
development, as in the data. Shocks are normal with a variance that rises with the agricultural
share. Shock sensitivity is
$\beta_d=-(1+0.5\,\text{agri}-0.1d)+T\,\rho_d(X)$, with resilience effect
$\rho_d(X)=0.8(1-d/10)\,\mathbf 1\{\text{internet}>0\}$. Capability therefore protects the poorer
deciles more, and only in connected countries. Income levels depend non-linearly on the
covariates and on $T$, so capability is also correlated with growth itself. Decile incomes are
log-normal with a spread that increases with inequality. We draw $N=150$ and $N=500$ countries,
set the budget at $k=0.3$ and run 200 replications. For every replication we compute the best
depth-2 tree under the true benefits, the oracle, and measure each method's regret against it
under the $\epsilon$-weighted welfare criterion.

\paragraph{Comparators.} EW-DRPL learns the tree from the equity-weighted orthogonal scores. The
utilitarian rule learns it from $\epsilon=0$ scores and is evaluated under $\epsilon$. The plug-in
rule uses $\hat m_{d,1}-\hat m_{d,0}$ without the doubly robust correction. The IPW-only rule drops
the outcome model. The raw-score rule uses the non-orthogonalised slope signal $\psi_d$ in place of
$\psi^*_d$. The naive rule ranks countries by the fitted interaction effects from a pooled
regression of $Y$ on $S$, $S\times T$ and $S\times T\times Z$ with decile fixed effects, the
standard specification in the applied literature.

\begin{table}[t]
\centering
\caption{Monte Carlo results}\label{tab:mc}
\scriptsize
\resizebox{\textwidth}{!}{% Monte Carlo summary for EW-DRPL. Requires \usepackage{booktabs}.
% Cells: mean (MC standard error) over 200 replications.
\begin{tabular}{lcccccc}
\toprule
 & \multicolumn{3}{c}{$N=150$} & \multicolumn{3}{c}{$N=500$} \\
\cmidrule(lr){2-4}\cmidrule(lr){5-7}
Method & $\varepsilon=0$ & $\varepsilon=1$ & $\varepsilon=2$ & $\varepsilon=0$ & $\varepsilon=1$ & $\varepsilon=2$ \\
\midrule
\multicolumn{7}{l}{\emph{Panel A: regret vs.\ oracle tree (true welfare)}} \\
EW-DRPL & 0.058 (0.002) & 0.070 (0.002) & 0.070 (0.005) & 0.028 (0.001) & 0.047 (0.002) & 0.051 (0.003) \\
EW-DRPL-crossfit & 0.057 (0.001) & 0.073 (0.002) & 0.079 (0.005) & 0.032 (0.001) & 0.050 (0.002) & 0.056 (0.003) \\
Utilitarian-DR & 0.058 (0.002) & 0.078 (0.002) & 0.106 (0.008) & 0.028 (0.001) & 0.052 (0.002) & 0.085 (0.007) \\
Plug-in & 0.046 (0.002) & 0.062 (0.002) & 0.069 (0.007) & 0.027 (0.001) & 0.043 (0.002) & 0.048 (0.004) \\
IPW-only & 0.064 (0.001) & 0.076 (0.002) & 0.074 (0.006) & 0.053 (0.001) & 0.075 (0.001) & 0.077 (0.005) \\
Raw-psi-DR & 0.066 (0.001) & 0.078 (0.002) & 0.083 (0.006) & 0.050 (0.001) & 0.069 (0.002) & 0.084 (0.006) \\
Naive-TWFE & 0.078 (0.001) & 0.105 (0.002) & 0.156 (0.009) & 0.077 (0.000) & 0.106 (0.001) & 0.166 (0.008) \\
Oracle & 0.000 (0.000) & 0.000 (0.000) & 0.000 (0.000) & 0.000 (0.000) & 0.000 (0.000) & 0.000 (0.000) \\
\midrule
\multicolumn{7}{l}{\emph{Panel B: share of oracle-treated countries recovered}} \\
EW-DRPL & 0.393 (0.010) & 0.376 (0.011) & 0.355 (0.009) & 0.584 (0.010) & 0.505 (0.013) & 0.419 (0.011) \\
EW-DRPL-crossfit & 0.405 (0.009) & 0.368 (0.009) & 0.355 (0.008) & 0.556 (0.008) & 0.488 (0.010) & 0.408 (0.008) \\
Utilitarian-DR & 0.393 (0.010) & 0.351 (0.010) & 0.336 (0.010) & 0.584 (0.010) & 0.467 (0.010) & 0.454 (0.010) \\
Plug-in & 0.478 (0.012) & 0.422 (0.011) & 0.396 (0.010) & 0.619 (0.012) & 0.532 (0.014) & 0.454 (0.013) \\
IPW-only & 0.363 (0.009) & 0.330 (0.010) & 0.321 (0.008) & 0.420 (0.008) & 0.332 (0.008) & 0.293 (0.008) \\
Raw-psi-DR & 0.349 (0.009) & 0.334 (0.009) & 0.326 (0.009) & 0.435 (0.009) & 0.373 (0.010) & 0.338 (0.009) \\
Naive-TWFE & 0.118 (0.006) & 0.325 (0.005) & 0.327 (0.004) & 0.093 (0.005) & 0.342 (0.003) & 0.338 (0.003) \\
Oracle & 1.000 (0.000) & 1.000 (0.000) & 1.000 (0.000) & 1.000 (0.000) & 1.000 (0.000) & 1.000 (0.000) \\
\midrule
\multicolumn{7}{l}{\emph{Panel C: ATE of $\rho$ (average over countries and deciles): bias (MC SE) and RMSE}} \\
ATE-DR & -0.057 (0.076) & 1.075 &  & 0.010 (0.005) & 0.072 &  \\
ATE-DR-raw & 0.052 (0.034) & 0.477 &  & -0.000 (0.006) & 0.086 &  \\
ATE-plugin & 0.036 (0.023) & 0.330 &  & 0.011 (0.005) & 0.069 &  \\
ATE-IPW & 0.077 (0.037) & 0.526 &  & 0.022 (0.006) & 0.084 &  \\
ATE-naive & 0.291 (0.006) & 0.303 &  & 0.300 (0.003) & 0.303 &  \\
\bottomrule
\end{tabular}
}
\begin{minipage}{\textwidth}\footnotesize\smallskip
\emph{Notes:} 200 replications; Monte Carlo standard errors in parentheses. Regret is the
difference between the welfare of the oracle depth-2 tree and that of the learned rule, both
evaluated with the true benefits under the indicated inequality aversion $\epsilon$. Budget
$k=0.3$. ``crossfit'' learns the tree on four country folds and applies it to the fifth.
Propensity scores are clipped to $[0.1,0.9]$.
\end{minipage}
\end{table}

\paragraph{Results.} Table~\ref{tab:mc} and Appendix Figure~\ref{fig:regret} report the results.
Four findings stand out.

First, the standard interaction regression is a poor targeting device. Its regret exceeds that of
EW-DRPL in every cell of the design, by between 30 per cent ($N=150$, $\epsilon=0$) and 71 per cent
($N=500$, $\epsilon=2$). At $\epsilon=0$ it recovers only about a tenth of the countries the oracle
would treat, fewer than random allocation would. Its implied average resilience effect is biased
upward by about 0.30, against a true average of about 0.18, because a pooled interaction model
cannot separate capability from the characteristics it is correlated with.

Second, equity weighting matters when the planner is strongly inequality averse. At $\epsilon=2$
the utilitarian rule's regret is 0.104 at $N=150$ and 0.084 at $N=500$, against 0.067 and 0.048 for
EW-DRPL, reductions of 35 and 43 per cent. At $\epsilon=1$ the reduction is smaller (9 and 15 per
cent), and at $\epsilon=0$ the two rules coincide by construction.

Third, orthogonalisation improves targeting and, in larger samples, estimation. Replacing
$\psi^*_d$ by the raw Robinson signal raises regret at $N=500$ by 0.019 to 0.033 and the RMSE of
the average resilience effect from 0.065 to 0.082. At $N=150$, however, the orthogonal score has
the larger RMSE (0.56 against 0.36), because the correction terms divide by estimated propensities
and shock variances. Dropping the outcome model (IPW-only) raises regret in every cell.

Fourth, the doubly robust correction does not improve targeting over a well-specified plug-in rule
in this design. The plug-in rule, which ranks countries by $\hat m_{d,1}-\hat m_{d,0}$, has lower
regret at $N=150$, the same regret at $N=500$, and a lower RMSE for the average effect at $N=150$ (0.33 against 0.56). Its
advantage is variance: it omits the inverse-weighted corrections. The doubly robust rule is
insurance against misspecification of the slope model, which our design does not impose. We
therefore report both in the application. Clipping the propensity score at $[0.1,0.9]$ rather than
$[0.05,0.95]$ cuts the RMSE of the doubly robust average effect at $N=150$ from 1.08 to 0.56 and
leaves regret essentially unchanged (Appendix Table~\ref{tab:mc_trim}).

The design is favourable to the estimator in several respects. Shocks are normal and independent
across countries, the response is linear in the shock, capability is redrawn for every spell, and
exposure and shock sizes are known. The application departs from each of these. The practical
lesson for cross-country work with around 150 countries is that average effects will be estimated
imprecisely even when targeting rules are informative, and that limited overlap between countries
with and without capability is the binding constraint.
